\section{integrali di superficie}

\begin{theorem}[significato geometrico del determinante jacobiano]
  \label{th:significato_geometrico_jacobiano}
  Sia $L\colon \RR^m \to \RR^n$ una funzione lineare
  e sia $V$ l'immagine di $L$.
  Sia $M\colon V \to \RR^m$ una isometria.
  Sia $E\subset \RR^m$.
  Allora si ha 
  \begin{equation}\label{eq:significato_geometrico_jacobiano}
    \abs{\det (M\circ L)}\cdot \abs{E} = \sqrt{\det (L^t L)}\cdot \abs{E}.
  \end{equation}
  Visto che le isometrie preservano le aree
  è naturale definire la $m$-area di $L(E)$ in $\RR^n$ 
  come la misura $m$-dimensionale di $M(L(E))$ in $\RR^m$.
  Per il teorema~\ref{th:significato_geometrico_determinante}
  la $m$-area di $L(E)$ è data dal lato sinistro di~\eqref{eq:significato_geometrico_jacobiano}
  che, guardando il lato destro, non dipende dalla scelta dell'isometria $M$:
  \[
    m\text{-}\area(L(E)) = \sqrt{\det (L^t L)}\cdot \abs{E}.
  \]
\end{theorem}
%
\begin{proof}
Estendiamo $M$ ad una isometria su tutto $\RR^n$, $\tilde M\colon \RR^n \to \RR^n$.
Dunque $\tilde M$ si rappresenta con una matrice ortogonale, cioè $\tilde M^t = \tilde M^{-1}$.
Sia $P\colon \RR^n \to \RR^m$ la proiezione ortogonale $P(x_1,\dots,x_n) = (x_1,\dots,x_m)$.
Dunque $M$ è la restrizione di $P\tilde M$ a $V$
e $M\circ L = P\tilde M L$.
Osserviamo che $P^t P\colon \RR^n\to \RR^n$ è l'applicazione 
lineare che azzera le ultime $n-m$ coordinate
ma visto che $M$ manda $V$ in $\RR^m$ si ha $P^t P \tilde ML = ML$. 
Dunque 
\[
 (\det(M\circ L))^2 
 = \det((P \tilde M L)^t (P \tilde M L)) 
 = \det(L^t \tilde M^t P^t P \tilde M L) 
 = \det(L^t \tilde M^t M L) = \det(L^t L)
\]
come volevamo dimostrare.
\end{proof}

\begin{definition}[integrale di superficie]
  \label{def:integrale_superficie}
Sia $B\subset \RR^k$ un aperto limitato misurabile e sia $A\subset \RR^n$.
Data $\phi\colon \bar B\subset \to A$ di classe $C^1$ 
e data $f\colon A\subset \RR^n\to \RR$ di classe $C^0$
definiamo
\begin{equation}\label{eq:def_integrale_superficie}
  \int_\phi f \, d\sigma 
  \defeq \int_B f(\phi(x)) \sqrt{\det ((D\phi)^t D\phi)}\, dx.
\end{equation}
Le ipotesi su $\phi$, e $f$ garantiscono 
che la funzione integranda al lato destro di~\eqref{eq:def_integrale_superficie} è
continua e dunque integrabile su $B$, con integrale finito. 
Per mettere in evidenza la variabile di integrazione si potrà anche scrivere 
\[
  \int_\phi f \, d\sigma = \int_\phi f(\phi(x)) \, d\sigma(x)
\]
intendendo quindi
$d\sigma(x) = \sqrt{\det ((D\phi)^t D\phi)}\, dx$.

Si può estendere la definizione al caso in cui $B$ non è limitato. 
In tal caso l'integrale potrebbe non esistere e potrebbe essere infinito.
\end{definition}

\begin{theorem}[invarianza per riparametrizzazioni]
Siano $B,C$ aperti misurabili limitati di $\RR^k$ e sia $A\subset \RR^n$.
Sia $\phi\colon \bar B \to A$ di classe $C^1$,
sia $f\colon A\to \RR$ di classe $C^0$ e sia 
$u\colon C \to B$ di classe $C^1$ 
con $u\colon C\to B$ bigettiva e $\det Du \neq 0$.
Allora 
\begin{align*}
  \int_{\phi\circ u} f \, d\sigma
  &= \int_{\phi} f \, d\sigma.
\end{align*}
\end{theorem}

\begin{proof}
Si ha
\begin{align*}
  \int_{\phi\circ u} f \, d\sigma 
  &= \int_C f(\phi(u(x))) \sqrt{\det ((D(\phi\circ u))^t D(\phi\circ u))}\, dx \\
  &= \int_C f(\phi(u(x))) \sqrt{\det ((Du)^t (D\phi)^t D\phi Du)}\, dx \\
  &= \int_C f(\phi(u(x))) \sqrt{\det ((D\phi)^t D\phi)}\cdot \sqrt{\det ((Du)^t Du)}\, dx \\
  &= \int_B f(\phi(y)) \sqrt{\det ((D\phi)^t D\phi)}\, dy\\
  &= \int_{\phi} f \, d\sigma
\end{align*}
avendo fatto il cambio di variabile:
\[
  y = u(x), \qquad dy = \abs{\det Du} dx = \sqrt{\det ((Du)^t Du)}\, dx.
\]
\end{proof}

Se $D\phi(x)$ ha rango massimo, i vettori colonna 
$\frac{\partial \phi}{\partial x_1},\dots,\frac{\partial \phi}{\partial x_k}$
sono linearmente indipendenti e dunque 
generano uno spazio vettoriale $k$-dimensionale 
che si chiama spazio tangente alla superficie definita da $\phi$ nel punto $\phi(x)$.


\section{ipersuperfici}

Se $k=n-1$, $\phi\colon B\subset \RR^{n-1} \to A\subset \RR^n$ 
rappresenta una \emph{ipersuperficie} in $\RR^n$ o superficie di codimensione $1$.
Se $D\phi$ ha rango massimo in un punto $x$, 
i vettori $\frac{\partial \phi}{\partial x_1},\dots,\frac{\partial \phi}{\partial x_{n-1}}$ 
sono linearmente indipendenti e dunque
generano lo spazio tangente alla superficie definita da $\phi$ nel punto $\phi(x)$.
Per ogni $v\in \RR^n$ possiamo considerare il determinante $\det(v, D\phi)$ 
della matrice quadrata ottenuta aggiungendo la colonna $v$ alla matrice $D\phi$.
Chiamiamo $\xi_\phi$ il vettore che rappresenta l'applicazione lineare $v\mapsto \det(v, D\phi)$, 
cioè $\xi_\phi$ è definito univocamente dalla proprietà
\[
  \xi_\phi \cdot v = \det(v,D\phi).
\]
Chiaramente $\xi_\phi$ è ortogonale a tutte le colonne di $D\phi$ perché
$\xi_\phi \cdot \phi_{x_i} = \det(\phi_{x_i}, D\phi) = 0$,
dunque $\xi_\phi$ è un vettore ortogonale allo spazio tangente alla superficie.
Definiamo $\nu_\phi$ come il versore $\nu_\phi=\xi_\phi/\abs{\xi_\phi}$, 
che è uno dei due versori normali alla ipersuperficie.
La direzione di $\nu_\phi$ non cambia se si riparametrizza la superficie
e il suo verso non cambia se la riparametrizzazione ha jacobiano positivo.

Se $\phi$ non è iniettivo è possibile che la superficie si autointersechi 
o addirittura che si sovrapponga più volte su se stessa. 
In questo senso diremo che la superficie ha una \emph{molteplicità}.

Per quanto riguarda il modulo di $\xi_\phi$ si ha 
\begin{align*}
   \abs{\xi_\phi} &= \det (\nu_\phi, D\phi)
   = \sqrt{\det \enclose{(\nu_\phi, D\phi)^t (\nu_\phi,D\phi)}}  \\
   &= \sqrt{\det \begin{pmatrix}
     1 & 0 \\
     0 & (D\phi)^t D\phi
   \end{pmatrix}}
  = \sqrt{\det((D\phi)^t D\phi)}
\end{align*}
dunque 
\[
  \xi_\phi = \nu_\phi \sqrt{\det (D\phi)^t D\phi}
\]
e per ogni $v\in \RR^n$ si ha
\begin{equation}\label{eq:493444}
\det(v, D\phi)
= \xi_\phi \cdot v
= (\nu_\phi \cdot v) \sqrt{\det((D\phi)^t D\phi)}.
\end{equation}

Formalmente, se $\sigma$ rappresenta l'integrazione di superficie vista in precedenza, 
si ha
\[
  \xi_\phi(x) \, dx = \nu_\phi\, d\sigma.
\]

Nel caso particolare $n=3$, se $\phi(u,v) = (x(u,v), y(u,v), z(u,v))$ è la 
parametrizzazione di una superficie in $\RR^3$, si ha 
\[
  (\xi_\phi(x),v)  = \det\enclose{v, \frac{\partial \phi}{\partial u}, \frac{\partial \phi}{\partial v}}
\]
e quindi, in base alla definizione~\ref{def:prodotto_vettore},
risulta
\[
  \xi_\phi = \frac{\partial \phi}{\partial u} \times \frac{\partial \phi}{\partial v}
\]
da cui
\[
  \nu_\phi\, d\sigma = \enclose{\frac{\partial \phi}{\partial u} \times \frac{\partial \phi}{\partial v}}\, du\, dv.
\]

\section{Esempi}

\subsection{area della sfera}

Osservando che la funzione identità: $x\colon \RR^n\to \RR^n$ 
ha divergenza costante 
$\div x = \frac{\partial x_1}{\partial x_1} + \dots + \frac{\partial x_n}{\partial x_n} = n$,
e osservando che $x$ è normale alla sfera $\partial B_R$, 
si può usare il teorema della divergenza per calcolare l'area della sfera di raggio $R$:
\[
\omega_n R^n = \abs{B_R} 
  = \frac 1 n \int_{B_R} \div x
  = \frac 1 n \int_{\partial B_R} x \cdot \nu_{\partial B_R}\, d\sigma
  = \frac 1 n \int_{\partial B_R} R\, d\sigma
  = \frac{R}{n} \cdot \area(\partial B_R)
\]
da cui
\begin{equation}\label{eq:area_sfera}
  \area(\partial B_R) 
  = n \omega_n R^{n-1}.
\end{equation}

Lo stesso risultato si può ottenere direttamente. 
Se $\theta\colon B\to \RR^n$ è una qualunque parametrizzazione della sfera $\partial B_1$,
possiamo parametrizzare la palla $B_R$ con la funzione $\phi \colon B\times [0,R] \to \RR^n$ 
definita da $\phi(u,\rho) = \rho\theta(u)$.
Si osserva allora che $D\phi = (\rho D\theta, \theta)$ e, 
osservando che $\frac{\partial \theta}{\partial u_i}$ 
è tangente alla sfera, e quindi ortogonale a $\theta$, mentre $\abs{\theta}=1$, 
risulta
\[
  D\phi^t D\phi = \begin{pmatrix}
  \rho^2 D\theta^t D\theta & \rho D\theta^t \theta \\
  \rho \theta^t D\theta & 1
  \end{pmatrix} = \begin{pmatrix}
  \rho^2 D\theta^t D\theta & 0 \\
  0 & 1
  \end{pmatrix}
\]
da cui l'elemento di volume sulla palla è dato da 
\[
  dV = \rho^{n-1} \sqrt{\det(D\theta^t D\theta)}\, du\, d\rho
     = \rho^{n-1} d\sigma\, d\rho
\]
dove $d\sigma$ è l'elemento di area sulla sfera unitaria $\partial B_1$.
Allora si ha
\[
\abs{B_R} 
  = \int_{\partial B_1} \int_0^R \rho^{n-1}\, d\rho\, d\sigma 
  = \frac{R^n}{n} \int_{\partial B_1} d\sigma 
  = \frac{R^n}{n} \area(\partial B_1)
\]
da cui si ottiene nuovamente~\eqref{eq:area_sfera}.

Il calcolo precedente può essere svolto in maniera più 
informale usando con più disinvoltura le notazioni del calcolo 
differenziale. 
Se chiamiamo $x$ la variabile che descrive il punto nella palla $B_R$,
chiamiamo $\theta$ la variabile che descrive il punto sulla sfera unitaria
poniamo $\rho=\abs{x}$, $\theta = \frac{x}{\abs{x}}$ e dunque $x = \rho \theta$,
i passaggi visti sopra si possono esprimere più espressivamente 
scrivendo $dx = \rho^{n-1} d\theta \, d\rho$ dove $d\theta$ è 
l'elemento d'area nell'integrazione sulla sfera unitara.

\subsection{grafico di una funzione}

Se $B\subset \RR^n$ è misurabile e $f\colon B\to \RR$, possiamo 
rappresentare la ipersuperficie $y=f(x)$ in $\RR^{n+1}$
tramite la parametrizzazione $\phi\colon B\to \RR^{n+1}$, 
$\phi(x) = (x, f(x))$.
Si può allora osservare che 
\[
  D\phi = \begin{pmatrix}
    \id \\
    Df
  \end{pmatrix},
  \qquad
  (D\phi)^t D\phi = \id + (Df)^t Df = \id + \nabla f \otimes \nabla f
\]
dove $\id$ è la matrice identità $n\times n$ e 
$v \otimes w = v w^t$ è il prodotto tensore tra i vettori (colonna) 
$v$ e $w$ cioè la matrice con componenti $(v \otimes w)_{i,j} = v_i w_j$.

Per verifica diretta si osserva che 
$(\id + v \otimes v)$ è una matrice simmetrica che 
ha $v$ come autovettore con autovalore $1+\abs{v}^2$ 
e che ogni vettore ortogonale a $v$ è un autovettore con autovalore $1$.
Dunque il determinante di $(\id + v \otimes v)$ è dato dal prodotto degli autovalori,
e quindi
\[
  \sqrt{\det((D\phi)^t D\phi)} = \sqrt{1 + \abs{\nabla f}^2}.
\]
La formula per calcolare l'area del grafico di $f$ è dunque
\[
  \area(f) = \int_B \sqrt{1 + \abs{\nabla f}^2}\, dx.
\]
Non è difficile calcolare il versore normale al grafico della funzione 
$f$, basta infatti osservare che il vettore $v=(-D f, 1)$ 
è ortogonale ad ogni colonna $\frac{\partial \phi}{\partial x_i} = (e_i, 1)$ 
della matrice jacobiana $D\phi$ e dunque, rinormalizzando, 
si ottiene la normale al grafico di~$f$:
\[
  \nu_f = \frac{\begin{pmatrix}
  -\nabla f\\
  1
  \end{pmatrix}}{\sqrt{1 + \abs{\nabla f}^2}}.
\]
Dunque, se volessimo calcolare un flusso attraverso la superficie, possiamo usare 
la formula 
\[
  \nu_f\, d\sigma = \begin{pmatrix}
  -\nabla f\\
  1
  \end{pmatrix} dx.
\]

\begin{example}[area dell'iperboloide]
Consideriamo l'iperboloide $H$ di equazione $z = x^2 - y^2$
 con $x^2+y^2\le R$, grafico della funzione $f(x,y) = x^2 - y^2$ su 
 $B_R = \{(x,y)\in \RR^2\colon x^2+y^2\le R\}$.
 Si ha $\nabla f = (2x, -2y)$ e dunque, 
 usando le coordinate polari, 
 si ottiene
 \begin{align*}
 \area(H) 
  &= \int_{B_R} \sqrt{1+(\nabla f)^2}\, dx\, dy
  = \int_{B_R} \sqrt{1 + 4(x^2+y^2)}\, dx\, dy \\
  & = \int_0^{2\pi} \int_0^R \sqrt{1 + 4\rho^2}\cdot \rho\, d\rho\, d\theta
  = 2 \pi \Enclose{\frac{(1+4R^2)^{\frac 3 2}}{12}}_0^R
  = \frac{\pi}{6} \Enclose{\enclose{1+4R^2}^{\frac 3 2} - 1}.
 \end{align*}
\end{example}


\subsection{superfici di rotazione}


Se $\gamma\colon [a,b]\to \RR^2$ è una curva di classe $C^1$, 
posto $\gamma(t) = (\rho(t), z(t))$, 
possiamo rappresentare la superficie di rotazione 
$S$ di equazione $x^2+y^2 = \rho(t)^2$, $z=z(t)$ tramite la parametrizzazione
$\phi(\theta, t) = (\rho(t)\cos\theta, \rho(t)\sin\theta, z(t))$.
Si ha 
\begin{align*}
  D\phi &= \begin{pmatrix}
  -\rho(t)\sin \theta & \rho'(t)\cos\theta \\
  \rho(t)\cos\theta & \rho'(t)\sin\theta \\
  0 & z'(t) 
  \end{pmatrix},\\
  \sqrt{\det (D\phi)^t D\phi} 
    &= \sqrt{\det \begin{pmatrix}
    \rho(t)^2 & 0 \\
    0 & \rho'(t)^2 + z'(t)^2
    \end{pmatrix}}
    = \rho(t) \sqrt{\rho'(t)^2 + z'(t)^2}\\
    &= \rho(t) \,\abs{\gamma'(t)}.
\end{align*}
Dunque si ottiene la \emph{formula di Guldino}
\index{formula!di Guldino}%
\index{Guldino, formula di}%
\begin{equation}\label{eq:formula_guldino}
 \area(S) = \int_a^b \int_0^{2\pi} \rho(t) \abs{\gamma'(t)}\, d\theta\, dt 
  = 2\pi \int_a^b \rho(t) \abs{\gamma'(t)}\, dt
  = 2\pi \int_\gamma \rho\, ds.
\end{equation}
Se denotiamo con 
\[
 \bar \rho = \dashint_\gamma \rho = \frac{\int_\gamma \rho}{\ell(\gamma)}
\]
il baricentro della curva $\gamma$, si ottiene 
$\area(S) = 2\pi \bar \rho \ell(\gamma)$ cioè
l'area di $S$ coincide con l'area si un cilindro 
con base sulla curva $\gamma$ e altezza pari alla 
lunghezza della circonferenza percorsa dal baricentro 
della curva.

\begin{example}[area della superficie di un toro]
Consideriamo il toro ottenuto facendo ruotare una circonferenza di 
raggio $r$ attorno ad un asse che dista $R$ dal centro della circonferenza, con $R>r$.
Il baricentro di una circonferenza è il centro della stessa, 
dunque l'area del toro è $2\pi R \cdot 2\pi r = 4\pi^2 R r$.
\end{example}
